\SourcePage{606}{609}
\section{General theory of scattering}

In classical mechanics, the collision of two particles is specified completely
by their velocities and by the impact parameter (the distance at which they
would pass one another if there were no interaction). In quantum mechanics the
question itself must be posed differently: for motion at specified velocities,
a trajectory has no meaning, and consequently neither does an impact
parameter. The task of the theory is then only to calculate the probability
that, after the collision, the particles will be deflected---or, as one says,
\emph{scattered}---through a particular angle. Here we consider so-called
elastic collisions, in which the particles undergo no transformations and, if
they are composite, their internal states remain unchanged.

Like every two-body problem, an elastic collision reduces to the scattering of
one particle, with the reduced mass, by the field $U(r)$ of a fixed force
centre\footnote{We neglect the spin--orbit interaction of the particles (if
they possess spin). By assuming a centrally symmetric field, we also exclude
from consideration processes such as electron scattering by molecules.}. The
reduction is made by passing to the frame in which the centre of mass of the
two particles is at rest. Let the scattering angle in this frame be $\theta$.
It is connected by simple formulae with the deflection angles $\theta_1$ and
$\theta_2$ of the two particles in the ``laboratory'' frame, in which one of
the particles (the second) was at rest before the collision:
\begin{equation}
 \tan\theta_1=\frac{m_2\sin\theta}{m_1+m_2\cos\theta},\qquad
 \theta_2=\frac{\pi-\theta}{2},
 \tag{123.1}
\end{equation}
where $m_1$ and $m_2$ are the particle masses (see I, \S~17). In particular,
when the masses are equal ($m_1=m_2$), this gives simply
\begin{equation}
 \theta_1=\frac\theta2,\qquad \theta_2=\frac{\pi-\theta}{2};
 \tag{123.2}
\end{equation}

\SourcePage{607}{610}
the sum $\theta_1+\theta_2=\pi/2$, so that the particles fly apart at a right
angle.

Throughout the remainder of this chapter, unless expressly stated otherwise,
we use the centre-of-mass frame and denote the reduced mass of the colliding
particles by $m$.

A free particle moving along the positive $z$-axis is represented by a plane
wave, which we write as $\psi=e^{ikz}$; thus we choose a normalization for
which the current density in the wave equals the particle velocity $v$. Far
from the centre, the scattered particles are represented by an outgoing
spherical wave $f(\theta)e^{ikr}/r$, where $f(\theta)$ is a function of the
scattering angle $\theta$ (the angle between the $z$-axis and the direction of
the scattered particle). This function is called the \emph{scattering
amplitude}. Hence the exact wave function solving the Schrödinger equation
with potential energy $U(r)$ must have, at large distances, the asymptotic
form
\begin{equation}
 \psi\approx e^{ikz}+\frac{f(\theta)}r e^{ikr}.
 \tag{123.3}
\end{equation}

The probability per unit time for a scattered particle to cross the surface
element $dS=r^2d\Omega$ (where $d\Omega$ is the element of solid angle) is
$vr^{-2}|f|^2dS=v|f|^2d\Omega$\footnote{This argument tacitly assumes that the
incident particle beam is bounded by a broad, but finite, aperture---broad so
as to avoid diffraction effects---as in actual scattering experiments. There
is consequently no interference between the two terms in (123.3): the square
$|\psi|^2$ is evaluated at points where the incident wave is absent.}. Its
ratio to the current density of the incident wave is
\begin{equation}
 d\sigma=|f(\theta)|^2d\Omega.
 \tag{123.4}
\end{equation}
This quantity has the dimensions of area and is called the \emph{effective
cross-section} (or simply the \emph{cross-section}) for scattering into the
solid-angle element $d\Omega$. On putting
$d\Omega=2\pi\sin\theta\,d\theta$, we obtain the cross-section
\begin{equation}
 d\sigma=2\pi\sin\theta|f(\theta)|^2d\theta
 \tag{123.5}
\end{equation}
for scattering into the angular interval from $\theta$ to
$\theta+d\theta$.

The solution of the Schrödinger equation describing scattering in the central
field $U(r)$ must plainly be axially symmetric about the $z$-axis, the
direction of the incident particles. Any such solution may be represented as
a

\SourcePage{608}{611}
superposition of continuous-spectrum wave functions for particles moving in
the given field with specified energy $\hbar^2k^2/2m$, with orbital angular
momenta of various magnitudes $l$ and vanishing $z$-components (these
functions do not depend on the azimuthal angle $\varphi$ about the $z$-axis
and are therefore axially symmetric). The required wave function thus has the
form
\begin{equation}
 \psi=\sum_{l=0}^\infty A_lP_l(\cos\theta)R_{kl}(r),
 \tag{123.6}
\end{equation}
where the $A_l$ are constants and the radial functions $R_{kl}(r)$ satisfy
\begin{equation}
 \frac1{r^2}\frac d{dr}\left(r^2\frac{dR_{kl}}{dr}\right)
 +\left[k^2-\frac{l(l+1)}{r^2}-\frac{2m}{\hbar^2}U(r)\right]R_{kl}=0.
 \tag{123.7}
\end{equation}
The coefficients $A_l$ must be chosen so that (123.6) has the asymptotic form
(123.3) at large distances. We shall show that this requires
\begin{equation}
 A_l=\frac1{2k}(2l+1)i^l\exp(i\delta_l),
 \tag{123.8}
\end{equation}
where $\delta_l$ are the phase shifts of the functions $R_{kl}$. This will at
the same time express the scattering amplitude in terms of these phases.

The asymptotic form of $R_{kl}$ is given by (33.20):
\[
 \begin{split}
 R_{kl}&\approx\frac2r\sin\left(kr-\frac{l\pi}{2}+\delta_l\right)\\
 &=\frac1{ir}\{(-i)^l\exp[i(kr+\delta_l)]
 -i^l\exp[-i(kr+\delta_l)]\}.
 \end{split}
\]
Substitution of this expression and (123.8) into (123.6) gives the asymptotic
wave function
\begin{equation}
 \psi\approx\frac1{2irk}\sum_{l=0}^\infty(2l+1)P_l(\cos\theta)
 [(-1)^{l+1}e^{-ikr}+S_le^{ikr}],
 \tag{123.9}
\end{equation}
where we have introduced the notation
\begin{equation}
 S_l=\exp(2i\delta_l).
 \tag{123.10}
\end{equation}
On the other hand, after the same rearrangement, the expansion (34.2) of a
plane wave reads
\[
 e^{ikz}\approx\frac1{2ikr}\sum_{l=0}^\infty(2l+1)P_l(\cos\theta)
 [(-1)^{l+1}e^{-ikr}+e^{ikr}].
\]

\SourcePage{609}{612}
We see that all terms containing $e^{-ikr}$ cancel, as they must, in the
difference $\psi-e^{ikz}$. The coefficient of $e^{ikr}/r$ in this difference,
which is the scattering amplitude, is therefore
\begin{equation}
 f(\theta)=\frac1{2ik}\sum_{l=0}^\infty(2l+1)(S_l-1)P_l(\cos\theta).
 \tag{123.11}
\end{equation}
This formula expresses the scattering amplitude in terms of the phases
$\delta_l$ (H.~Faxen, J.~Holtsmark, 1927)\footnote{Of fundamental interest is
the inverse problem of recovering the scattering potential when the phases
$\delta_l$ are presumed known. This problem was solved by
\emph{I.~M.~Gelfand, B.~M.~Levitan, and V.~A.~Marchenko}. In principle, to
determine $U(r)$ it is sufficient to know $\delta_0(k)$ as a function of the
wave vector over the whole range $0\leq k<\infty$, together with the
coefficients $a_n$ in the asymptotic expressions (as $r\to\infty$)
\[
 R_{n0}\approx a_ne^{-\kappa_nr}/r
 \quad(\kappa_n=\sqrt{2m|E_n|}/\hbar)
\]
for the wave functions of states belonging to discrete (negative) energy
levels $E_n$, if such levels exist. Determining $U(r)$ from these data reduces
to solving a certain linear integral equation. A systematic account may be
found in: V.~de~Alfaro and T.~Regge, \emph{Potential Scattering}. Moscow: Mir,
1966.}.

Integrating $d\sigma$ over all angles gives the total scattering
cross-section $\sigma$, which is the ratio of the total probability per unit
time that a particle is scattered to the current density in the incident
wave. Substituting (123.11) in
\[
 \sigma=2\pi\int_0^\pi|f(\theta)|^2\sin\theta\,d\theta
\]
and using the mutual orthogonality of Legendre polynomials with different
$l$, together with
\[
 \int_0^\pi P_l^2(\cos\theta)\sin\theta\,d\theta=\frac2{2l+1},
\]
we obtain the following expression for the total cross-section:
\begin{equation}
 \sigma=\frac{4\pi}{k^2}\sum_{l=0}^\infty(2l+1)\sin^2\delta_l.
 \tag{123.12}
\end{equation}

Each term of this sum is the \emph{partial cross-section} $\sigma_l$ for the
scattering of particles with the specified orbital angular

\SourcePage{610}{613}
momentum $l$. The greatest possible value of this cross-section is
\begin{equation}
 \sigma_{l\,\max}=\frac{4\pi}{k^2}(2l+1).
 \tag{123.13}
\end{equation}
Comparison with (34.5) shows that the number of particles scattered with
angular momentum $l$ can be four times the number of such particles in the
incident flux. This circumstance is a purely quantum effect associated with
interference between scattered and unscattered particles.

We shall also find it convenient below to use the \emph{partial scattering
amplitudes} $f_l$, defined as the coefficients in the expansion
\begin{equation}
 f(\theta)=\sum_{l=0}^\infty(2l+1)f_lP_l(\cos\theta).
 \tag{123.14}
\end{equation}
According to (123.11), they are related to the phases $\delta_l$ by
\begin{equation}
 f_l=\frac1{2ik}(S_l-1)=\frac1{2ik}(e^{2i\delta_l}-1),
 \tag{123.15}
\end{equation}
and the partial cross-sections are
\begin{equation}
 \sigma_l=4\pi(2l+1)|f_l|^2.
 \tag{123.16}
\end{equation}

\ProblemHeading

Express the scattering amplitude in terms of phase shifts in two dimensions.
The field is $U=U(\rho)$, where $\rho=\sqrt{x^2+z^2}$, and the particle flux
is directed along the $z$-axis.

\SolutionHeading In two dimensions, far from the scatterer the wave function
is a superposition of a plane wave and an outgoing cylindrical wave:
\begin{equation}
 \psi=e^{ikz}+f(\varphi)\frac{e^{ik\rho}}{\sqrt{-i\rho}}.
 \tag{1}
\end{equation}
Here $\varphi$ is the angle between the $z$-axis and the scattering direction,
and $f(\varphi)$ is the scattering amplitude, whose dimension in two
dimensions is the square root of length. The factor
$-i=\exp(-i\pi/2)$ beneath the radical has been introduced to simplify the
subsequent formulae. The scattering cross-section per unit length along the
$y$-axis is
\[
 d\sigma=|f|^2d\varphi.
\]
It has the dimension of length.

The wave function must be expanded in functions with a definite projection
$m$ of angular momentum on the $y$-axis, of the form
$Q_m(\rho)e^{im\varphi}$. At large distances from the scatterer, the radial
functions differ from the free-motion functions obtained in the problem to
\S~34 only by a phase shift:
\[
 Q_m(\rho)\approx i^m\sqrt{\frac2{\pi k\rho}}
 \sin\left[k\rho-\frac\pi2\left(m-\frac12\right)+\delta_m\right],
\]

\SourcePage{611}{614}
with $\delta_m=\delta_{-m}$. Repeating the argument of this section and using
the plane-wave expansion from the problem to \S~34, we find that the function
with asymptotic form (1) is given by the series
\[
 \psi=\sum_{m=-\infty}^{\infty}e^{i\delta_m}Q_m(\rho)e^{im\varphi},
\]
and that the scattering amplitude is
\begin{equation}
 f(\varphi)=\frac1{i\sqrt{2\pi k}}
 \sum_{m=-\infty}^{\infty}(e^{2i\delta_m}-1)e^{im\varphi}.
 \tag{2}
\end{equation}
Integration gives the total cross-section
\[
 \sigma=\int_0^{2\pi}|f|^2d\varphi
 =\sum_{m=-\infty}^{\infty}\sigma_m,
 \qquad \text{where}\qquad
 \sigma_m=\sigma_{-m}=\frac4k\sin^2\delta_m.
\]
It is readily verified that
\begin{equation}
 \operatorname{Im}f(0)=\sqrt{\frac{k}{8\pi}}\,\sigma,
 \tag{3}
\end{equation}
which is the optical theorem in two dimensions (see formula (125.9) below).
